\chapter{Diskussion}
\label{ch:06_diskussion}

Dieses Kapitel führt die theoretischen Anforderungen und architektonischen Grenzen der untersuchten
Rendering-Architekturen aus Kapitel~\ref{ch:04_konzeption_realisierung} sowie die praktischen Messergebnisse aus
Kapitel~\ref{ch:05_evaluation} zusammen und stellt die Ergebnisse in Zusammenhang mit den in
Kapitel~\ref{sec:0102_goal_and_core_question} formulierten Forschungsfragen, um diese zu beantworten.
Daraus resultierend wird das finale heuristische Entscheidungsmodell abgeleitet und anhand der erarbeiteten Ergebnisse
in Form von Umsetzungsempfehlungen dargelegt.
Abschließend werden die in dieser Arbeit gesammelten Erkenntnisse, getroffenen Entscheidungen und umgesetzten Maßnahmen
einer kritischen Reflexion unterzogen.


\section{Zusammenführung und Interpretation der Ergebnisse}
\label{sec:0601_interpretation_results}

% Einleitung.
Die in Kapitel ~\ref{sec:0102_goal_and_core_question} dargelegten Forschungsfragen lassen sich für die in dieser Arbeit
gewählten Rendering-Pipelines, deren Implementation die absoluten Basis-Architekturen der jeweiligen
Rendering-Paradigmen darstellen, individuell beantworten.

% Forward Basis Performance.
Wie in Kapitel ~\ref{ch:03_acedemic_and_industry_state} dargelegt und
Kapitel~\ref{subsec:050101_evaluation_tagets_and_hypothesis} postuliert, wird das klassische
Single-Pass-Forward-Rendering auf Performance-Ebene theoretisch primär durch zwei Flaschenhälse definiert:
Fragment-Overdraw und die Verarbeitung großer Anzahlen von Lichtquellen.
Wie die Ergebnisse aus~\ref{subsec:050203_overdraw_eval} und~\ref{subsec:050403_shading_overdraw_eval} darstellen,
wird auf Hardwareebene moderner \ac{GPU}s bereits ein extrem großer Teil an Overdraw-Engpässen mitigiert.
Auch wenn Overdraw, wie~\textcite[Abschn.~1,~Punkt~5]{Thaler2011DeferredRendering} darlegt, nie gänzlich eliminiert
werden kann, sind die inhärenten Funktionen moderner \acp{GPU} bereits ausreichend, um auch ohne zusätzliche
Algorithmen, wie ein dedizierter Z-Prepass, Overdraw als primären Bottleneck weitestgehend zu umgehen.
Versuch~\ref{fig:5_4_1_light_count_means} bestätigt hingegen die grundlegende These des signifikanten Bottlenecks
hinsichtlich extrem hoher Lichtquellen.
Weiterhin stellt Versuch~\ref{subsec:050202_obejct_count_eval} den wenig evaluierten Bottleneck auf \ac{CPU}-Ebene für
Forward-Pipelines in Form von ansteigende \ac{CPU}-Dispatch-Zeiten durch State-Changes, Culling und
Lichtquellen-Suchläufe.
Wie in Abbildung~\ref{fig:5_2_2_cpu_scaling_logic_vs_render} dargestellt, ist dieser Effekt allerdings vergleichsweise
weniger signifikant für die Performance des Gesamtsystems.

% Forward Stilisierung.
Die Beantwortung der grundlegenden Frage nach der Umsetzung von Forward-Pipelines bezüglich selektiver Stilisierung
wirft, wie in Kapitel~\ref{subsec:040201_forward_pipeline} beschrieben, eine grundlegende Restriktion auf.
Forward-Pipelines weisen zwar in der reinen Verwaltung verschiedener Shader und Materialien signifikante praktische
Vorteile gegenüber Deferred-Paradigmen auf~\cite[Abschn.~9.2]{Chapter9Deferred}, aufgrund ihrer grundlegenden
Eigenschaft, Shading-Operationen auf Objektebene (Object Space) darzustellen, sind sie jedoch nicht in der Lage,
klassische Bildraum (Screen-Space) basierte Verfahren, wie etwa Kernel-basierte Effekte, selektiv darzustellen.
Wie in Kapitel~\ref{subsec:040201_forward_pipeline} aufgezeigt, ist es je nach Anwendungsfall jedoch möglich, dass
vergleichbare \ac{NPR}-Effekte auf Objektebene existieren, die ersatzweise verwendet werden können (beispielsweise die
Verwendung von Inverted-Hull statt Sobel-Outlines).

% Single Pass.
Eine weitere Einschränkung der untersuchten Forward-Pipeline abseits von Performance und selektiver Stilisierungen ist
die begrenzte Anzahl maximal darstellbarer Lichtquellen durch \textit{Uniform-Buffer-Limits}.
Diese Begrenzung lässt sich allerdings weniger auf die Natur von Forward-Rendering-Architekturen im Allgemeinen
zurückführen und mehr auf die Natur von Single-Pass-Architekturen~\cite[Abschn.~1]{Thaler2011DeferredRendering}.
Diese Einschränkung trifft also ebenfalls auf die untersuchte Deferred-Uber-Pipeline zu.
In der Deferred-Uber-Pipeline ist durch den Singe-Pass Umstand zudem der Zwang zu \textit{Pre-Lit}-Verwahren, wie in
Kapitel~\ref{subsec:040202_deferred_uber_pipeline} aufgezeigt, festgestellt worden.
Soll ein bestimmter \ac{NPR}-Effekt also unbedingt selektiv als reiner Post-Lit-Effekt verarbeitet werden, um die
integrität eines Post-Processing Effektes zu gewährleisten (vgl.\ ~\ref{fig:040202_pre-post-lit_comparison}) muss auch
innerhalb der Deferred-Paradigmen auf Multi-Pass-Verfahren zurückgegriffen werden.

% Deferred Vorteil.
Deferred Rendering in seinen zwei ausgeführten Formen verspricht, wie in Kapitel~\ref{ch:03_acedemic_and_industry_state}
erörtert, diverse Vorteile gegenüber klassischem Forward-Rendering und für beide untersuchte Pipelines konnte die
Grundannahme der effizienteren Verarbeitung hoher Anzahlen von Lichtquellen bestätigt werden
(vgl.\ ~Kapitel~\ref{fig:5_4_1_light_count_means}).

% Deferred Drawbacks.
In Bezug auf selektive Stilisierung in Deferred Pipelines und Effizient bei hoher Szeneentropie konnte zudem auf
Performanceebene die Hypothese der Warp-Divergence
(vgl.\ ~Kapitel~\ref{subsec:020203_warp_divergence_and_register_pressure}) für die Deferred-Uber-Pipeline in
Versuch~\ref{subsec:050501_entropy_eval} bestätigt werden, auch wenn eine eindeutige
Zuweisung des Phänomens auf Hardwareebene durch den empirischen Aufbau der Versuche nicht endgültig verifiziert werden
kann (etwa welchen anteiligen Einfluss Registerdruck hier hat).
Obwohl der Deferred-Multi-Pass-Ansatz der Deferred-Volume-Pipeline in diesem Kontext keine Schwächen gezeigt hat,
konnten in den Versuchen~\ref{subsec:050302_bandwidth_eval} und~\ref{subsec:050402_light_intensity_and_overlap_eval} die
durch Blending in die Höhe getriebenen \ac{RMW}-Zyklen eindeutig identifiziert werden
(vgl.\ Kapitel~\ref{subsec:020202_memory_bandwidth}).
Die Deferred-Multi-Pass-Ansätze zeigen zudem, wie in Kapitel~\ref{subsec:040203_deferred_volume_pipeline} dargelegt, den
weiteren Nachteil der erhöhten \ac{VRAM}-Belastung durch größere G-Buffer-Formate, um Artefakte wie Banding (siehe
Abbildung~\ref{fig:hydra_banding_comparison}) zu vermeiden.

% Kapitel 3 Vergleich.
Es ist anzumerken, das viele der Benannten Bottlenecks und Drawbacks der einzelnen Paradigmen, wie in
Kapitel~\ref{sec:0302_industry_state} und~\ref{sec:0301_academic_state} dargelegt, zumindest teilweise durch moderne
Verfahren und hybride Ansätze kompensiert werden können.
Verfahren wie beispielsweise \textit{Clustered Forward}~\cite{olssonClusteredDeferredForward2012} oder
\textit{Forward+}~\cite{haradaForwardBringingDeferred2012} erweitern die
Basisparadigmen hierfür Grundlegend.
Moderne Entwicklungen aus der Spieleindustrie zeigen hingegen die Idee auf, hybride Verfahren mit einer Kombination aus
Forward und Deferred Ansätzen für den Ausgleich der jeweiligen Schwächen einer Rendering-Architektur zu
nutzen~\cite{muhammadPrettyFramesPRAGMATA2026, gameworksToonRenderingHiFi, valientDeferredRenderingKillzone}.


\section{Heuristisches Entscheidungsmodell}
\label{sec:0602_interpretation_results}

Dieses Kapitel konsolidiert die bisherigen Erkenntnisse über die Basis-Paradigmen der untersuchten Rendering-Pipelines
in einem heuristischen Entscheidungsmodell.
Wie in Kapitel~\ref{ch:03_acedemic_and_industry_state} dargelegt, ist die Wahl einer Rendering-Architektur in der Praxis
stark vom angestrebten Zielbild und dessen spezifischen Anforderungen
abhängig.
Die nachfolgende Auswertung differenziert zwischen deterministischen Ausschlusskriterien (\textbf{Hard Limits}), welche
durch architektonische oder \ac{API}-bedingte Restriktionen definiert sind
(vgl.\ Kapitel~\ref{ch:04_konzeption_realisierung}), und empirischen Skalierungsheuristiken (\textbf{Soft Limits}), die
Restriktionen und Kipppunkte auf Hardwareebene beschreiben (vgl.\ Kapitel~\ref{ch:05_evaluation}).
Die Synthese dieser Ergebnisse wird in Abbildung~\ref{fig:heuristic_decision_model} visualisiert.

\paragraph{Deterministische Ausschlusskriterien (Hard Limits)}
Im Folgenden werden die architektonischen und Grafik-\ac{API}-bedingten Restriktionen evaluiert, die eine Pipeline-Wahl
deterministisch ausschließen:

\begin{description}
    \item[Selektive Stilisierung von Bildraum-Effekten] ~\\
    Durch die inhärente Limitierung von Forward-Rendering-Paradigmen Geometrien ausschließlich auf Objekten
    (Object-Space), ist Forward-Rendering in seiner Grundform ungeeignet, um Bildraumeffekte (Screen-Space)
    wie Kuwahara-Filter oder Sobel-Outlines selektiv darzustellen (vgl.\ ~\ref{subsec:040201_forward_pipeline}).
    Es muss auf Deferred-Paradigmen oder visuell äquivalente Effekte (vgl.\ Abb.~\ref{fig:040201_outline_comparison})
    zurückgegriffen werden.

    \item[Licht-Quantität (1)] ~\\
    Wird die Anzahl der aktiven Lichtquellen in einer Szene erwartet größer als das von der Grafik-\ac{API} zugrunde
    gelegte Limit für Uniform-Buffer (etwa 512 in OpenGL), so ist der Einsatz einer Multi-Pass-Pipeline mit aufgeteilten
    Lichtberechnungspässen erforderlich (vgl.\ ~\ref{subsec:040201_forward_pipeline}).
    Simple Single-Pass Architekturen wie eine naive Forward-Rendering-Pipeline oder Deferred-Pipelines mit einem
    einzigen monolithischen Uber-Shader zur Lichtberechnung, sind hier nicht hinreichend
    (vgl.\ Abb.~\ref{fig:5_4_1_light_count_means}).

    \item[NPR-Integrität (Pre-Lit vs Post-Lit)] ~\\
    Bei der Anwendung von kernelbasierten Effekten, wie beispielsweise dem Kuwahara-Filter, in selektiven
    Deferred-Pipelines erweisen sich Single-Pass-Deferred-Pipelines als unzureichend, sofern die vollständige Integrität
    des resultierenden Effekts gewahrt werden muss (vgl.\ ~\ref{subsec:040202_deferred_uber_pipeline}).
    Aufgrund der signifikanten Zunahme des algorithmischen Aufwands, die Beleuchtungsberechnungen
    der notwendigen Umgebungspixel für jeden bearbeiteten Bildpunkt durchzuführen, ist es erforderlich die
    Lichteinwirkung bereits vor der Anwendung des eigentlichen \ac{NPR}-Effekts ine einem vorgelagerten Pass zu
    berechnen (vgl.\ ~\ref{subsec:040203_deferred_volume_pipeline}).
\end{description}

\paragraph{Empirische Skalierungsheuristiken (Soft Limits)}
Im Folgenden werden die graduellen Performance-Kipppunkte auf Hardwareebene erläutert, ab denen die Performance bei
zunehmender Szenenkomplexität unverhältnismäßig stark sinkt.

\begin{description}
    \item[Licht-Quantität (2)] ~\\
    Forward-Rendering ist hinreichend geeignet Szenen mit einer niedrigen Anzahl an Lichtquellen darzustellen, in
    Umgebungen mit einer signifikanten Anzahl von Lichtquellen ist von der Verwendung von Forward-Paradigmen zur
    Lichtauswertung allerdings abzusehen (vgl.\ Abb.~\ref{fig:5_4_1_light_count_means}).
    Der Forward-inhärente quadratische Aufwand von $\mathcal{O}(N_{\text{obj}} \cdot N_{\text{light}})$ ist dem linearen Aufwand von
    $\mathcal{O}(N_{\text{obj}} + N_{\text{light}})$ in Deferred signifikant unterlegen
    (vgl.\ ~\ref{subsec:020101_forward_rendering} u.~\ref{subsec:020102_deferred_rendering}).

    \item[Monolithische Shader-Strukturen] ~\\
    Wenn Szenen mit einer extrem hohen Stil- oder Material-Entropie vorliegen, ist bei der Verwendung von monolithischen
    Shader-Strukturen mit einer großen Anzahl von \texttt{if/else}- oder \texttt{switch}-Bedingungen, die sich auf nicht
    triviale Shading-Operationen beziehen, in der selektiven Stilisierung mit Einschränkungen zu rechnen
    (vgl.\ ~\ref{subsec:050501_entropy_eval}).
    Dies gilt insbesondere für Deferred-Rendering-Paradigmen, die eine höhere Tendenz zu solchen monolithischen
    Shader-Strukturen durch vereinheitlichte Shading-Pässe aufweisen (vgl.\ ~\ref{subsec:020102_deferred_rendering}).

    \item[Light-Volumes] ~\\
    Bei der Verwendung von Light-Volumes zur Lichtberechnung, die im Deferred Rendering weit verbreitet sind, ist in
    komplexen Szenen mit einer hohen Lichtdichte und der daraus resultierenden großen Menge an Überlappungen der
    Lichteinflüsse mit Performanceeinbußen zu rechnen
    (vgl.\ ~\ref{subsec:050302_bandwidth_eval} u.~\ref{subsec:050402_light_intensity_and_overlap_eval}).
    Die durch Additives Blending der Lichteinflüsse erhöhten \ac{RMW}-Zyklen auf der \ac{GPU} können sich als
    signifikanter Performance-Engpass herausstellen (vgl.\ ~\ref{subsec:020202_memory_bandwidth}).
    In diesem Zusammenhang muss weiterhin beachtet werden, dass das Deferred Rendering grundlegend höhere
    Speicherbelastungen aufweist, um Artefaktbildungen zu vermeiden und damit eine höhere Grundlast auf der \ac{GPU}
    verursacht (vgl.\ ~\ref{subsec:040203_deferred_volume_pipeline}).
    Dies könnte bei kleinen und sehr speichertechnisch begrenzten Systemen einen signifikanten Faktor darstellen
    (vgl.\ ~\ref{subsec:050302_bandwidth_eval} u. ~\ref{sec:0506_macro_eval}).

    \item[Objekt-Quantität] ~\\
    In komplexen Szenen mit einer hohen Anzahl individueller Objekte weisen klassische Forward-Rendering-Pipelines eine
    erhöhte Last auf der Host-Seite (\ac{CPU}) durch insgesamt höhere Dispatch-Zeiten auf
    (vgl.\ ~\ref{subsec:050202_obejct_count_eval}).
    In Umgebungen mit wenig Rechenkapazität auf \ac{CPU}-Seite sind Deferred-Architekturen somit möglicherweise
    vorzuziehen, auch wenn dieser Effekt im Gesamtsystem eher marginal auffällt
    (vgl.\ Abb.~\ref{fig:5_2_2_cpu_scaling_logic_vs_render}).
\end{description}

\begin{figure}[htbp]
    \centering
    \resizebox{\textwidth}{!}{%
        \begin{tikzpicture}[
            >=Latex,
            font=\sffamily\small,
            base/.style={draw, align=center, minimum height=0.7cm, minimum width=4.5cm},
            startstop/.style={base, rounded corners, fill=gray!15},
            process/.style={base, fill=white},
            decision/.style={draw, diamond, aspect=2.6, inner sep=2pt, align=center, fill=white,
            font=\sffamily\footnotesize},
            branchact/.style={base, rounded corners, fill=gray!15, minimum width=4.5cm, minimum height=1.0cm},
            hardedge/.style={->, draw=red!75!black, thick, rounded corners=4pt},
            softedge/.style={->, draw=blue!75!black, thick, rounded corners=4pt},
            normedge/.style={->, draw=black, thick, rounded corners=4pt}
        ]

            % --- Legende ---
            \node[draw, fill=white, align=left, font=\sffamily\footnotesize, rounded corners, anchor=north west] at
            (-9.0, 1.0) {
                \textbf{Entscheidungspfade:} \\
                \textcolor{red!75!black}{\rule{0.6cm}{1.5pt}} Deterministische Hard Limits \\
                \textcolor{blue!75!black}{\rule{0.6cm}{1.5pt}} Empirische Soft Limits
            };

            % --- Knotenpunkte ---

            % Start
            \node[startstop] (init) at (0, 0) {Szenen- \& Anforderungsanalyse \\ \texttt{EvaluateHeuristics()}};

            % Hard Limits
            \node[decision] (d_api) at (0, -3.5) {\textbf{Uniform-Limit} \\ \textbf{überschritten?}};
            \node[decision] (d_npr) at (0, -7.5) {\textbf{Selektive Bildraum-} \\ \textbf{Effekte nötig?}};
            \node[decision] (d_prelit) at (6.5, -7.5) {\textbf{Strikte NPR-} \\ \textbf{Integrität gefordert?}};

            % Soft Limits
            \node[decision] (d_light) at (0, -11.5) {\textbf{Sehr hohe} \\ \textbf{Lichtquellen-Anzahl?}};
            \node[decision] (d_obj) at (0, -15.5) {\textbf{Sehr hohe Objektanzahl} \\
            \textbf{\& niedrige CPU-Ressourcen?}};
            \node[decision] (d_entropy) at (6.5, -20.5) {\textbf{Hohe Material-} \\ \textbf{Entropie?}};
            \node[decision] (d_vram) at (6.5, -24.5) {\textbf{Extreme Lichtdichte} \\ \textbf{oder VRAM-Limit?}};

            % Pipelines (Ziele)
            \node[branchact] (p_fwd) at (-6.5, -29.5) {\textbf{Forward Pipeline} \\
            \footnotesize{Single-Pass Object-Space}};
            \node[branchact] (p_uber) at (0, -29.5) {\textbf{Deferred-Uber Pipeline} \\
            \footnotesize{Single-Pass Screen-Space}};
            \node[branchact] (p_vol) at (10.5, -29.5) {\textbf{Deferred-Volume Pipeline} \\
            \footnotesize{Multi-Pass Screen-Space}};


            % --- Kontrollfluss ---

            % Start
            \draw[normedge] (init) -- (d_api);

            % API Limit (Hard)
            \draw[hardedge] (d_api.south) -- node[right, font=\footnotesize] {Nein} (d_npr.north);
            % Drop-Linie ganz rechts
            \draw[hardedge] (d_api.east) -| node[pos=0.15, above, font=\footnotesize] {Ja} ([xshift=2.0cm]p_vol.north);

            % Screen-Space NPR (Hard)
            \draw[hardedge] (d_npr.south) -- node[right, font=\footnotesize] {Nein} (d_light.north);
            \draw[hardedge] (d_npr.east) -- node[above, font=\footnotesize] {Ja} (d_prelit.west);

            % Pre-Lit Integrity (Hard)
            \draw[hardedge] (d_prelit.south) -- node[right, font=\footnotesize] {Nein} (d_entropy.north);
            % Drop-Linie rechts (X=11.5)
            \draw[hardedge] (d_prelit.east) -| node[pos=0.15, above, font=\footnotesize] {Ja}
            ([xshift=1.0cm]p_vol.north);

            % Soft Limit: Lichtquellen
            \draw[softedge] (d_light.south) -- node[right, font=\footnotesize] {Nein} (d_obj.north);
            % Sub-Routing über Hilfskoordinate zur Entropie-Abfrage
            \draw[softedge] (d_light.east) -- node[above, font=\footnotesize] {Ja} (4.5, -11.5) -- (4.5, -18.5) -|
            ([xshift=-1.5cm]d_entropy.north);

            % Soft Limit 2: Objekte & CPU
            \draw[softedge] (d_obj.west) -| node[pos=0.15, above, font=\footnotesize] {Nein} (p_fwd.north);
            % Zusammenführung mit dem 'Ja'-Pfad der Lichtquellen bei X=4.5
            \draw[softedge] (d_obj.east) -- node[above, font=\footnotesize] {Ja} (4.5, -15.5);

            % Material Entropy (Soft)
            \draw[softedge] (d_entropy.south) -- node[right, font=\footnotesize] {Nein} (d_vram.north);
            % Drop-Linie mittig rechts
            \draw[softedge] (d_entropy.east) -| node[pos=0.15, above, font=\footnotesize] {Ja}
            ([xshift=-0.5cm]p_vol.north);

            % VRAM / Overlap (Soft)
            % Drop-Linie zentral auf Uber-Pipeline
            \draw[softedge] (d_vram.west) -| node[pos=0.15, above, font=\footnotesize] {Ja} (p_uber.north);
            % Drop-Linie links von den roten Pfaden
            \draw[softedge] (d_vram.south) -- node[right, font=\footnotesize] {Nein} (6.5, -27.0) -|
            ([xshift=-1.5cm]p_vol.north);

        \end{tikzpicture}%
    }
    \caption[Heuristisches Entscheidungsmodell der Rendering-Architekturen]{UML-Aktivitätsdiagramm des heuristischen
    Entscheidungsmodells zur Wahl der optimalen Basis-Rendering-Architektur.
    Eine farbliche Trennung verdeutlicht die Gewichtung zwischen zwingenden Architekturvorgaben (Rot) und
    Hardware-Skalierungsgrenzen (Blau).}
    \label{fig:heuristic_decision_model}
\end{figure}


\section{Kritische Reflexion und Limitationen}
\label{sec:0603_reflektion_and_limits}

% Einleitung.
In der nachfolgenden Analyse werden die Grenzen des erstellten Evaluations-Frameworks \ac{HyDra}, die Lücken der
empirischen Evaluationen sowie sonstige limitierende Faktoren und Entscheidungen der vorliegenden Arbeit reflektiert.

% Fehlende Rendering-Paradigmen.
Eine Einschränkung des erstellten Frameworks ist im grundlegenden Umfang der implementierten Rendering-Verfahren zu
verzeichnen.
Die Beschränkung der Umgebung auf die gewählten drei Rendering-Architekturen ist im Sinne einer Evaluation der
Basisparadigmen zwar vertretbar, in der Praxis sind jedoch, wie in Kapitel ~\ref{sec:0302_industry_state} aufgezeigt,
wesentlich fragmentiertere Umgebungen vorhanden und die entsprechenden Basisparadigmen kommen kaum in ihrer Reinform
vor.
Was in der Praxis als \textquote{Standard} oder \textquote{Basis} moderner Rendering-Architekturen gilt, ist
entsprechend nur schwer einzugrenzen und zu evaluieren.
Dies zeigen mittlerweile weit verbreitete Ansätze wie \textit{Forward+}~\cite{haradaForwardBringingDeferred2012},
\textit{Clustered}-Ansätze~\cite{olssonClusteredDeferredForward2012},
\textit{Tiled Shading}~\cite{olssonTiledClusteredForward2012} oder die nutzung von
\textit{Visibility-Buffern} in Deferred-Architekturen~\cite{burnsVisibilityBufferCacheFriendly2013}
als Erweiterungen der Basisparadigmen.
Insbesondere im Hinblick auf die Forward-Rendering-Baseline besteht eine signifikante Lücke bezüglich der
Multi-Pass-Variante des Renderingverfahrens.

% Materiallimits.
Bestimmte grundlegende Materialeigenschaften, wie etwa Transluzenzen oder Spiegelungen, stellen ebenfalls gängige und
grundlegende Anforderungen an Echtzeiterendering-Systeme
dar~\cites[Abschn.~\textit{Deferred G-buffer}]{muhammadPrettyFramesPRAGMATA2026}
[Abschn.~\textit{G-Buffer : Our approach}]{valientDeferredRenderingKillzone}
[S.~887]{akenine-mollerRealtimeRendering2019}.
In der vorliegenden Untersuchung werden sie jedoch nicht berücksichtigt.
Dies erschwert eine ganzheitliche Evaluation und Ableitung für diverse Praxisfälle.

% Lichtevaluationen.
In der Industriepraxis gehören darüber hinaus verschiedene Typen von Lichtquellen, wie Richtungs- (Directional) oder
Flächenlichter (Area Lights), und die damit verbundenen Anforderungen, insbesondere das \textit{Shadow Mapping}, zum
Standardrepertoire einer Szene in der
Spieleentwicklung~\cites[Abschn.~2.4.1]{fonsecaDeferredShadingTutorial}
[Kap.~11.1.3.3]{gregoryGameEngineArchitecture2018}.
Das \ac{HyDra}-Framework weist hier ebenfalls Lücken durch die reine Verwendung von dynamisch generierten
Punktlichtquellen auf.

% Empirie.
Innerhalb der angewandten Evaluationsmetriken und des grundlegenden Ansatzes der empirischen Untersuchung können
ebenfalls signifikante Limitationen der implementierten Software sowie der daraus resultierenden Untersuchungen
festgestellt werden.
Im Rahmen der Analyse spezifischer, hardwareinterner Phänomene und Flaschenhälse erweisen sich empirische
Vorgehensweisen, die sich ausschließlich auf die Messung von Verarbeitungszeiten konzentrieren, häufig als unzureichend.
Ein umfassendes Bild über die Ursachen und Limitierungen der betreffenden Phänomene und Hardware-Kipppunkte ist so nur
schwer zu gewinnen.
Dies zeigen vor allem die Versuche ~\ref{subsec:050302_bandwidth_eval} und ~\ref{subsec:050501_entropy_eval}.
Zwar ermöglichen rein empirische Messungen eine Untersuchung des grundlegenden Verhaltens einer Rendering-Pipeline, sie
sind jedoch nicht dazu geeignet, entsprechende detaillierte Handlungsempfehlungen bezüglich Hardwareoptimierungen
abzuleiten.
Das detaillierte Auslesen von hardwarinternen Daten müsste durch die zusätzliche Integration von spezifischen
Hardware-Profilern wie \textit{NVIDIA Nsight} oder der \textit{Radeon Developer Tool Suite (RDTS)}
erfolgen~\cite[Kap.~18.1]{akenine-mollerRealtimeRendering2019}.

% Hardwarespezifika.
Über die strukturellen Lücken der vorliegenden Software hinaus ist anzumerken, dass die empirischen Evaluationen aus
Kapitel~\ref{ch:05_evaluation} in ihrem Geltungsbereich an die zugrunde liegende Hardware
(vgl.\ ~\ref{tab:hardware_specs}) gebunden sind.
Um beispielsweise die Hardware-Kipppunkte auf mobilen Geräten genauer bestimmen zu können, müsste die zugrundeliegende
Hardware ebenfalls als dynamischer Vektor dargestellt werden, um eine allgemeingültigere Evaluation zu ermöglichen.


\section{Lessons Learned}
\label{sec:0604_lessons_learned}

% Architekturentscheidungen.
Im Entwicklungsverlauf ist insbesondere die Verwendung von \textit{OpenGL}~\cite{KhronosOpenGLRegistry} sowie vor allem
der Bibliothek \textit{raylib}~\cites{Raylib, RaylibExamples} zur Abstraktion der OpenGL-\ac{API} als positiv zu
bewerten.
Ein signifikanter Anteil des rudimentären OpenGL-Overheads konnte durch die Verwendung von raylib vermieden werden.
Für die granulare Manipulation diverser Teile des Frameworks, wie etwa der direkten Manipulation der \ac{G-Buffer},
musste allerdings auch innerhalb von raylib auf Low-Level-Boilerplate-Code zurückgegriffen werden
(\texttt{rlgl}~\cite{raysan5RaylibSrcRlglh}).
Die Granularität der Software stellte sich in der Analyse als signifikante Herausforderung heraus, da die reine Menge an
Boilerplate-Code zur Ausführung der Rendering-Pipelines einen wesentlichen Beitrag zur Komplexität und damit zur
Fehleranfälligkeit der Software leistet.
Zukünftige Erweiterungen und Anpassungen der Software sollten daher streng modular ablaufen, um monolithische
Codestrukturen, vor allem hinsichtlich der Grafik-\ac{API}-Instruktionen und damit assoziierten C++-Code, zu vermeiden.

% Szenenaufbau & Tooling.
Im Rahmen diverser Implementationsschritte war, durch diesen granularen Ansatz und die daraus resultierende
Code-Komplexität, teilweise eine vollständige Neuentwicklung großer Teile des Codes erforderlich.
Der Wechsel von einer statischen, hardcodierten Umgebung der Rendering-Pipelines zu einem für die Hardware-Benchmarks
notwendigen Input-getriebenen Design stellte einen extremen Paradigmenwechsel dar.
Für alle potenziellen zukünftigen Untersuchungsvektoren muss demnach im Voraus eine Mindestflexibilität des Frameworks
eingehalten werden.

% Irrwege der Versuchsvektoren.
Für die Erzeugung belastbarer Ergebnisse der durchgeführten Evaluationen stellte die genaue Definition der
Untersuchungsvektoren und des daraus resultierenden Szeneaufbaus eine Herausforderung dar.
Auch in ihren grundlegendsten Formen existieren für die untersuchten Rendering-Architekturen eine extrem schwer
eingrenzbare Anzahl an Eigenschaften.
Die wichtigsten Eigenschaften innerhalb des extrem weit gefassten Themas des Echtzeit-Rendering zu isolieren, ohne sich
in irrelevante Methodiken oder unnötige Optimierungen zu verlieren, erforderte eine iterative und zeitintensive
Abwägung.
Auf Implementationsebene manifestierte sich dies insbesondere in der Erzeugung einer belastbaren Szene.
Es ist essenziell, dass die Szene selbst nicht zu synthetisch aufgebaut ist, da dies die Übertragbarkeit auf
reale Szenarien in der Industriepraxis einschränkt.
Gleichzeitig ist darauf zu achten, dass die Komplexität der Szene nicht zu hoch ist, um eine Kontrolle über die
untersuchten Variablen zu gewährleisten.
Es ist festzustellen, dass sich Metriken generell nur schwer isolieren lassen.
Eine Änderung der Beleuchtung beeinflusst etwa unweigerlich die Speicherbandbreite.
Die strikte Isolation und eindimensionale Gestaltung von Testvektoren innerhalb kompletter
Rendering-Architekturen kann daher weitestgehend als methodischer Trugschluss betrachtet werden.

% Telemetrie und Benchmark-Overhead.
Die entwickelten Benchmark- und Telemetrie-Module sind besonders bezüglich ihres resultierenden Aufwands hervorzuheben.
Der Entwicklungsaufwand eines deterministischen und belastbaren Messsystems
(vgl.\ ~\ref{sec:0404_telemetry_and_benchmark_orchestrator}) mit sauberen Hardwareschleifen (Ring-Buffer-Verfahren) hat
sich, neben der Implementation der Rendering- und Shading-Verfahren, als eine der größten Herausforderungen
herausgestellt.
In Zukunft ist von einer Entwicklung der Messgrundlagen im Nachgang der Kernentwicklungen abzusehen.
Telemetrie und Benchmark-Module sollten als integraler Bestandteil der Basisarchitektur konzipiert werden, nicht als
nachgelagerte Evaluationsebene.

% Benchmark Evaluation.
In Anlehnung an die schwierige Abgrenzung auf Implementationsebene wurde auch bei der Definition der genauen
Testszenarien in Kapitel ~\ref{ch:05_evaluation} die Frage sinnvoller Rahmenbedingungen aufgeworfen.
Aufgrund des Umfangs der untersuchten Thematik erweist sich die Entwicklung eines Designs für belastbare Ergebnisse von
Testszenarien und -variablen als anspruchsvoll.
Die darauffolgende Interpretation der Ergebnisse auf rein empirischer Basis erforderte zudem eine weitreichende und
zeitintensive Deduktion anhand der vorliegenden Literatur.
Ohne tiefgreifende Profiler (bspw.\ \textit{NVIDIA Nsight}) bleibt die Interpretation von Frame-Times und Bandbreiten
stark deduktiv und spekulativ.
Dies sollte in zukünftigen Erweiterungen der Software durch die Integration jener Profiler vermieden werden.
